\documentclass[11pt]{article}
\usepackage[T1]{fontenc}
\usepackage{lmodern}
\usepackage{amsmath,amssymb,amsthm,mathtools}
\usepackage{graphicx}
\usepackage[margin=1in]{geometry}
\usepackage{microtype}
\usepackage{enumitem}
\usepackage[hidelinks]{hyperref}

\newtheorem{theorem}{Theorem}
\newtheorem{lemma}[theorem]{Lemma}
\newtheorem{corollary}[theorem]{Corollary}
\newcommand{\A}{\mathcal A}
\newcommand{\D}{\mathcal D}
\newcommand{\Prim}{\operatorname{Prim}}
\newcommand{\hull}{\operatorname{hull}}
\newcommand{\Ad}{\operatorname{Ad}}

\title{Finite and Clopen Primitive Sets Are Sets of Synthesis\\
for Locally Elliptic Weighted Orlicz Algebras}
\author{Automated Mathematics Run \texttt{fa\_banach\_001}, lane 2}
\date{21 July 2026}

\begin{document}
\maketitle

\begin{abstract}
Question 8.12 of Carter's paper asks for geometric conditions on subsets of
the primitive and coadjoint-orbit spectra which imply spectral synthesis for
Hermitian weighted Orlicz algebras on unipotent $p$-adic groups.  We prove two
such conditions.  Every clopen primitive subset is synthetic for every
locally elliptic group covered by Carter's hypotheses.  Every finite closed
set of admissible primitive points is also synthetic; for a $p$-adic
unipotent group, CCR makes every primitive point admissible.  Consequently,
every finite or clopen subset of the $p$-adic coadjoint-orbit space is a set
of synthesis.  The clopen proof uses central idempotents in compact-open
corners and approximates them by exact finite-Hecke Riesz projections in
Carter's minimal ideal $j(C)$.  This cannot extend to all compact closed
sets: Malliavin's theorem already supplies compact nonsynthesis sets for the
unweighted algebra on the abelian group $U_2(k)\cong(k,+)$.
\end{abstract}

\paragraph{Status.}
Candidate substantial partial result, likely valid, subject to human review.
It supplies the broad geometric criteria ``finite'' and ``clopen'' for the
full nonabelian class in Question 8.12.  Infinite closed sets with nonempty
boundary remain unclassified, and compactness alone is provably insufficient.

\section{Source question}

The source is Max Carter, \emph{Weighted Orlicz $*$-algebras on locally
elliptic groups}, arXiv:2503.20735v2, Question 8.12 on page 29
\cite{Carter}.  For a $p$-adic unipotent matrix group $N$, Carter identifies
\[
 \Prim_*(L^\Phi(N,\omega))\cong
 \mathfrak n^*/\Ad^*(N)
\]
and asks which geometric properties of $C$ imply $\ker(C)=j(C)$.

\begin{figure}[ht]
  \centering
  \includegraphics[width=\textwidth]{figures/open_problem_crop.png}
  \caption{Question 8.12, rendered from page 29 of
  arXiv:2503.20735v2.}
\end{figure}

\section{Definitions and setup}

Let $G$ be locally elliptic, let $(\Phi,\Psi)$ be a complementary pair of
Young functions with $\Phi\in\Delta_2$, and let $\omega$ be an
$L^\Phi$-weight with $1/\omega\in L^\Psi(G)$.  Put
\[
  \A=L^\Phi(G,\omega)
\]
and assume that $\A$ is Hermitian.  Carter identifies irreducible
$*$-representations of $\A$ with irreducible unitary representations of $G$.
For closed $C\subseteq\Prim_*(\A)$ he defines
\[
 \ker(C)=\bigcap_{P\in C}P
\]
and, using a smooth functional calculus, constructs a closed two-sided ideal
$j(C)$ with $j(C)\subseteq\ker(C)$ and $\hull_*(j(C))=C$
\cite[Theorem 7.1]{Carter}.  The set $C$ is a \emph{set of synthesis} if
$j(C)=\ker(C)$.

Write $\D(G)$ for the compactly supported locally constant functions.  For a
compact-open subgroup $H\leq G$, put
\[
 e_H=\mu(H)^{-1}1_H.
\]
Then $e_H=e_H^*=e_H*e_H$, and $\pi(e_H)$ is projection onto the $H$-fixed
vectors of a unitary representation $\pi$.

\section{Proof intuition}

Every algebraic test function belongs to a finite-dimensional Hecke algebra.
If it vanishes under the representations in $C$, then so does its central
support projection.  Carter's smooth functional calculus recovers that
projection exactly, placing the test function in $j(C)$.

For a clopen $C$, fix $H$ and pass to the unital corner $e_H\A e_H$.  The
clopen split of its primitive spectrum produces a central self-adjoint
idempotent $p_H$ which is zero on $C$ and one on the complementary primitive
points.  The idempotent need not itself be a test function.  Approximate it
in the Orlicz norm by self-adjoint Hecke functions and take their Riesz
projections for the spectral cluster near one.  Those projections remain in
finite Hecke algebras, vanish exactly on $C$, and converge to $p_H$; hence
$p_H\in j(C)$.  Every $H$-smoothed element of $\ker(C)$ is multiplied
identically by $p_H$, so compact-open smoothing completes the proof.

For a finite family of admissible representations, the same smoothing makes
the evaluation range finite dimensional.  A bounded linear lift corrects
dense Hecke approximants to vanish exactly at every chosen point.

\section{The algebraic test-function lemma}

\begin{lemma}[Dense locally finite test algebra]\label{lem:test}
$\D(G)$ is dense in $\A$, and every $d\in\D(G)$ belongs to a
finite-dimensional convolution $*$-subalgebra of $\D(G)$.
\end{lemma}

\begin{proof}
The $\Delta_2$ hypothesis gives density of compactly supported functions in
the weighted Orlicz space.  A locally elliptic group has a compact-open basis,
so uniform approximation on compact sets by locally constant step functions,
together with local boundedness of $\omega$, gives density of $\D(G)$.

The compact support of $d$ lies in a compact-open subgroup $K$.  Uniform local
constancy on $K$ gives an open subgroup under whose left and right
translations $d$ is invariant.  Replacing it by its normal core gives an open
normal $L\triangleleft K$.  The $L$-bi-invariant functions supported in $K$
form a finite-dimensional convolution $*$-algebra, isomorphic after Haar
normalization to $\mathbb C[K/L]$, and contain $d$.
\end{proof}

\begin{lemma}[Algebraic kernel elements lie in $j(C)$]\label{lem:algebraic}
For every closed $C\subseteq\Prim_*(\A)$,
\[
 \D(G)\cap\ker(C)\subseteq j(C).
\]
\end{lemma}

\begin{proof}
Put $d\in\D(G)\cap\ker(C)$ in a finite-dimensional convolution
$*$-algebra $B$ using Lemma~\ref{lem:test}.  In its faithful regular
representation, $B$ is a finite-dimensional $C^*$-algebra.  Let $p_d$ be the
central support projection of $d$, so $p_d*d=d$.

For every $\pi$ whose kernel lies in $C$, the ideal
$\ker(\pi|_B)$ contains $d$, hence contains $p_d$.  Choose
$0<a<2\pi$ with $a\lVert p_d\rVert_1\leq1$.  Carter's Proposition 5.5 gives
$\varphi\in A_\gamma$ which is zero near $0$ and satisfies $\varphi(a)=1$.
Then $\varphi\{ap_d\}\in m(C)$, while the defining power series and
$p_d*p_d=p_d$ give
\[
 \varphi\{ap_d\}=\varphi(a)p_d=p_d.
\]
Thus $p_d\in j(C)$ and $d=p_d*d\in j(C)$.
\end{proof}

\section{Clopen synthesis}

\begin{lemma}[The clopen corner idempotent]\label{lem:corner}
Let $C\subseteq\Prim_*(\A)$ be clopen and $H\leq G$ compact-open.  In the
corner
\[
 \A_H=e_H*\A*e_H
\]
there is a central self-adjoint idempotent $p_H$ such that, for every
irreducible $\pi$ with $\pi(e_H)\ne0$,
\[
 \pi(p_H)=
 \begin{cases}
 0,&\ker\pi\in C,\\
 \pi(e_H),&\ker\pi\notin C.
 \end{cases}
\]
\end{lemma}

\begin{proof}
The corner is unital with identity $e_H$.  It is Hermitian: if
$x=x^*\in\A_H$, then $x+(1-e_H)$ is self-adjoint in the unitization of $\A$;
compressing its nonreal resolvents shows that $x$ has real spectrum in the
corner.  It is semisimple because $\A\subseteq L^1(G)$ embeds faithfully in
the group $C^*$-algebra, and so does the corner.

The standard idempotent-corner correspondence, obtained by compressing
irreducible representations to $\pi(e_H)\mathcal H_\pi$, identifies
$\Prim_*(\A_H)$ with
\[
 U_H=\{\ker\pi\in\Prim_*(\A):\pi(e_H)\ne0\}.
\]
Hence $Y=C\cap U_H$ and $U_H\setminus C$ are complementary closed subsets of
$\Prim_*(\A_H)$.  Put
\[
 I=\ker_{\A_H}(Y),\qquad
 J=\ker_{\A_H}(U_H\setminus C).
\]
Semisimplicity gives $I\cap J=0$.  If $I+J$ were proper, a maximal left ideal
containing it would produce a simple $\A_H$-module.  Hermitianness makes that
module unitarizable, so its primitive kernel would belong to both $Y$ and its
complement, a contradiction.  Thus $I+J=\A_H$.

Write $e_H=p_H+q_H$ with $p_H\in I$ and $q_H\in J$.  Since products between
$I$ and $J$ lie in their zero intersection, $p_H$ and $q_H$ are complementary
central idempotents.  The ideals are $*$-closed, and uniqueness of the
decomposition gives $p_H=p_H^*$.  Its represented values are as stated.
\end{proof}

\begin{lemma}[The corner idempotent belongs to $j(C)$]\label{lem:riesz}
For $C$ and $p_H$ as above, $p_H\in j(C)$.
\end{lemma}

\begin{proof}
Let $\D_H=e_H*\D(G)*e_H$.  It is dense in $\A_H$.  Choose self-adjoint
$d_n\in\D_H$ with $d_n\to p_H$ in the $\A$-norm.  Every $d_n$ belongs to a
finite-dimensional Hecke $C^*$-algebra $B_n\subseteq\D_H$ with identity
$e_H$.

Because $p_H$ is central, $d_n$ commutes with $p_H$, and the spectra of
$p_Hd_n$ and $(e_H-p_H)d_n$ lie respectively near $1$ and $0$.  For large
$n$ these clusters are separated by the circle $\Gamma$ centered at $1$ with
radius $1/3$.  The spectrum of $d_n$ in $B_n$ equals its spectrum in
$\A_H$: if an eigenvalue in the finite-dimensional algebra were absent from
the corner spectrum, the corner inverse would annihilate its nonzero spectral
projection.  Consequently the Riesz projection
\[
 r_n=\frac{1}{2\pi i}\int_\Gamma(ze_H-d_n)^{-1}\,dz
\]
belongs to $B_n\subseteq\D(G)$.  Continuity of Riesz projections gives
$r_n\to p_H$ in the $\A$-norm.

If $\ker\pi\in C$ and $\pi(e_H)\ne0$, then
$\pi(d_n)\to\pi(p_H)=0$, so the Riesz projection of $\pi(d_n)$ around
$1$ is zero for large $n$.  If $\pi(e_H)=0$, every corner element is already
annihilated.  Hence $r_n\in\D(G)\cap\ker(C)$.  Lemma~\ref{lem:algebraic}
gives $r_n\in j(C)$, and closedness gives $p_H\in j(C)$.
\end{proof}

\begin{theorem}[Clopen synthesis]\label{thm:clopen}
Every clopen subset of $\Prim_*(\A)$ is a set of synthesis.
\end{theorem}

\begin{proof}
Let $f\in\ker(C)$ and fix compact-open $H$.  Put
$f_H=e_H*f*e_H$.  Lemma~\ref{lem:corner} and separation by irreducible
representations give
\[
 f_H=p_H*f_H.
\]
Indeed both sides vanish on $C$; off $C$, $p_H$ is the identity in every
nonzero corner representation.  Lemma~\ref{lem:riesz} now gives
$f_H\in j(C)$.

As $H$ shrinks to the identity, Carter's compact-neighbourhood approximate
identity \cite[Proposition 2.16(ii)]{Carter} yields $f_H\to f$ in $\A$.
Therefore $f\in j(C)$.  Since Carter already gives
$j(C)\subseteq\ker(C)$, equality follows.
\end{proof}

\section{Finite admissible synthesis}

Call an irreducible $\pi$ \emph{admissible} if
$\dim\mathcal H_\pi^H<\infty$ for every compact-open $H$.

\begin{theorem}[Finite admissible synthesis]\label{thm:finite}
Every finite closed subset of $\Prim_*(\A)$ consisting of admissible points
is a set of synthesis.
\end{theorem}

\begin{proof}
Let $C=\{\ker\pi_1,\ldots,\ker\pi_r\}$ and $I=\ker(C)$.  For fixed $H$,
evaluation maps $\A_H=e_H\A e_H$ into the finite-dimensional algebra
\[
 \bigoplus_{i=1}^r\mathcal B(\mathcal H_{\pi_i}^H).
\]
The dense subspace $\D_H$ has the same image because a linear subspace of a
finite-dimensional space is closed.  Choose a bounded linear right inverse
$R_H$ from that image into $\D_H$.

For $f\in I$, approximate $f_H=e_H*f*e_H$ by $d_n\in\D_H$ and set
\[
 h_n=d_n-R_H\bigl((\pi_1(d_n),\ldots,\pi_r(d_n))\bigr).
\]
Then $h_n\in\D(G)\cap I$ and $h_n\to f_H$.  Lemma~\ref{lem:algebraic}
gives $f_H\in j(C)$, and the compact-open approximate identity gives
$f\in j(C)$.
\end{proof}

\begin{corollary}[$p$-adic unipotent orbit spaces]\label{cor:padic}
Let $N$ be a group of unipotent matrices over a $p$-adic field and let
$\A=L^\Phi(N,\omega)$ satisfy the hypotheses above.  Every finite or clopen
subset of
\[
 \mathfrak n^*/\Ad^*(N)\cong\Prim_*(\A)
\]
is a set of synthesis.
\end{corollary}

\begin{proof}
Carter's Theorem 8.10 states that $N$ is CCR and identifies its unitary dual
with the coadjoint-orbit space.  For compact-open $H$, the projection
$\pi(e_H)$ lies in $\pi(C^*(N))=\mathcal K(\mathcal H_\pi)$, hence has finite
rank.  Thus every irreducible is admissible.  CCR also makes primitive ideals
maximal, so finite subsets are closed.  Apply Theorems~\ref{thm:clopen} and
\ref{thm:finite}, together with Carter's Corollary 8.11.
\end{proof}

\section{Sharpness: compact closed sets can fail synthesis}

\begin{theorem}[Malliavin obstruction]\label{thm:malliavin}
Let $k$ be a non-archimedean local field and identify
$U_2(k)\cong(k,+)$.  For the unweighted Hermitian algebra
$\A=L^1(k)$, there is a compact closed set
\[
 C\subseteq \Prim_*(\A)\cong\widehat{k}\cong k
\]
which is not a set of synthesis.  Consequently, neither closedness nor
compactness alone can replace finiteness or clopenness in
Corollary~\ref{cor:padic}.
\end{theorem}

\begin{proof}
The case $\Phi(t)=|t|$ and $\omega=1$ gives $\A=L^1(k)$ and lies in the
source framework.  Malliavin proved that spectral synthesis fails for
$L^1(G)$ whenever the locally compact abelian group $G$ is noncompact
\cite{Malliavin}.  In its compact-set formulation, every nondiscrete locally
compact abelian dual contains a compact subset on which synthesis fails
\cite[p.~358]{deLeeuwMirkil}.  The additive group of $k$ is noncompact and
Pontryagin self-dual, so this gives a compact closed
$C\subseteq\widehat{k}$ and a closed ideal $I$ satisfying
\[
 \hull(I)=C,\qquad I\subsetneq\ker(C).
\]
Carter's Theorem 7.1 makes $j(C)$ the minimal closed ideal for this hull, so
$j(C)\subseteq I\subsetneq\ker(C)$.  Hence $j(C)\ne\ker(C)$.
\end{proof}

\section{Verification, limitations, and novelty}

No numerical computation is used.  The proof audit checks the Hermitian
corner, the complementary-kernel central idempotent, equality of spectra in
the finite Hecke algebra and the corner, Riesz-projection convergence, and
the fact that no uniform finite-corner lift is needed as $H$ shrinks.

The theorem does not classify infinite closed subsets with nonempty boundary.
Theorem~\ref{thm:malliavin} shows that this is a genuine obstruction rather
than merely a limitation of the proof: even compact closed subsets can fail
synthesis in the abelian $p$-adic case.
In a non-Hausdorff primitive spectrum, ``compact-open'' alone need not imply
closed; Carter's $j(C)$ is defined for closed $C$, while the clopen theorem
has exactly the required closedness.  Local indexes and bounded arXiv searches
on 21 July 2026 combined the exact source id with ``spectral synthesis'',
``$p$-adic unipotent'', ``clopen'', ``compact-open'', ``finite sets'',
``central idempotent'', and ``weighted Orlicz''.  No exact theorem was found;
novelty confidence is moderate pending expert review.

\paragraph{Human-review recommendation.}
Verify Lemma~\ref{lem:corner}'s primitive corner correspondence and
Lemma~\ref{lem:riesz}'s equality of the finite-Hecke and corner spectra.
Those are the decisive new structural steps.

\begin{thebibliography}{9}
\bibitem{Carter}
M.~Carter,
\emph{Weighted Orlicz $*$-algebras on locally elliptic groups},
Studia Math. \textbf{287} (2026), no.~2, 145--180;
arXiv:2503.20735v2.

\bibitem{Malliavin}
P.~Malliavin,
\emph{Impossibilit\'e de la synth\`ese spectrale sur les groupes ab\'eliens
non compacts},
Publ. Math. IH\'ES \textbf{2} (1959), 61--68.

\bibitem{deLeeuwMirkil}
K.~de Leeuw and H.~Mirkil,
\emph{Translation-invariant function algebras on abelian groups},
Bull. Soc. Math. France \textbf{88} (1960), 345--370.
\end{thebibliography}

\end{document}
